\writeSectionFrame{\presentationTitle}{Fazit}
\begin{frame}{Anwendungsfälle}
	\begin{itemize}
 		\item Es ist wichtig, dass bestimmte Schritte in einer Reihenfolge ausgeführt werden
 		\item Klient muss nicht wissen, welches Objekt seine Anfrage bearbeitet
 		\item Länge der Kette ist nicht im Vornherein klar
 		\item Beispiel: Plugin-Mechanismus
 	\end{itemize}
\end{frame}
 
\note{
	Dass eine Anfrage eine Kette von Objekten entlang weitergeleitet werden muss,
	kommt recht häufig vor. Dieses Muster ist dann sinnvoll, wenn der Klient nicht
	wissen muss, welches Objekt seine Anfrage bearbeitet, oder wenn es sogar
	passieren kann, dass gar kein Objekt die Anfrage bearbeitet. Ketten können auch immer
	bis zum Ende durchgeführt werden, unabhängig davon, welche Bearbeiter bereits zum Zuge
	kamen.	Auch die Länge der
	Kette muss nicht von Anfang an klar sein, sondern kann während der Laufzeit bpsw. durch 
	Konfigurationen	bestimmt werden. Plugins, die nachgeladen werden, können sich in
	eine Kette einreihen, wenn sie Funktionalität nachrüsten sollen. 
}

\realworld{Java EE / Jakarta EE}{Servlet Filters}
\note{
    Java Servlet API verwendet das Chain-of-Responsibility-Muster für Servlet Filters. Servlet Filters sind eine Möglichkeit, HTTP-Anfragen und -Antworten in einer Filterkette zu durchlaufen. Jeder Filter in der Kette kann die Anfrage modifizieren oder sie an den nächste Filter weiterleiten.
}

\realworld{Apache Commons Chain}{Implementierung des Chain-of-Responsibility-Musters in Java}
\note{
    Das Apache Commons Chain-Framework ist eine Implementierung des Chain-of-Responsibility-Musters in Java. Es ermöglicht das Erstellen von Verarbeitungsketten, bei denen jeder Schritt (Handler) in der Kette entscheiden kann, ob er die Anfrage verarbeiten oder an den nächsten Schritt weiterleiten soll.
}

\realworld{Spring Security}{Sicherheitsüberprüfungen durch eine Filter Chain}
\note{
    Spring Security verwendet das Chain-of-Responsibility-Muster, um Sicherheitsüberprüfungen durch eine Filter Chain zu handhaben. Jede Sicherheitsüberprüfung wird als separater Filter implementiert, z.B. Authentifizierung, Autorisierung, CORS, CSRF-Schutz. Diese Filter sind verkettet, sodass jede Sicherheitsüberprüfung die Anfrage entweder verarbeiten oder weiterleiten kann.
}

\realworld{Log4j Logging Framework}{Logging-Events}
\note{
    Log4j verwendet das Chain-of-Responsibility-Muster für die Handhabung von Logging Events. In Log4j gibt es eine Kette von Appendern (z.B. ConsoleAppender, FileAppender), die nacheinander durchlaufen werden, um zu bestimmen, wie ein Log-Ereignis verarbeitet werden soll. Jeder Appender kann entscheiden, ob er das Ereignis verarbeitet oder an den nächsten Appender weiterleitet.
}

\realworld{JavaFX}{Ereignisverarbeitung}
\note{
    JavaFX verwendet das Chain of Responsibility-Muster in der Ereignisverarbeitung, indem es eine Kette von Event Filters und Handlers verwendet. 
}

\vorteil{Flexibilität bei der Handhabung von Anfragen}
\note{
    Das Muster ermöglicht es, neue Handler einfach hinzuzufügen, ohne den bestehenden Code ändern zu müssen. Neue Handler können in die Kette eingefügt werden, und die Reihenfolge kann angepasst werden.
}

\vorteil{Wiederverwendbarkeit von Handlern}
\note{
    Einzelne Handler können in verschiedenen Ketten wiederverwendet werden. Jeder Handler ist unabhängig von den anderen und kann in unterschiedlichen Kombinationen verwendet werden.
}

\vorteil{Entkopplung von Sender und Empfänger}
\note{
    Der Sender einer Anfrage (oder eines Ereignisses) muss nicht wissen, welcher Handler die Anfrage tatsächlich bearbeitet. Die Handlers sind lose gekoppelt, was zu einer besseren Modularität und Flexibilität führt.
}

\vorteil{Einfache Erweiterbarkeit}
\note{
    Neue Verarbeitungsschritte oder Logik können leicht hinzugefügt werden, indem neue Handler erstellt und der Kette hinzugefügt werden. Bestehende Handler müssen nicht geändert werden.
}

\vorteil{Reihenfolge der Aufrufe kann kontrolliert werden}
\note{
    Durch die Verkettung der Handler kann die Reihenfolge der Ausführung einfach kontrolliert und gesteuert werden.
}

\vorteil{Single-Responsibility-Prinzip}
\note{
    Jeder Handler hat eine klar definierte Aufgabe.
}

\nachteil{Komplexität der Kettenverwaltung}
\note{
    Das Verwalten und Debuggen von komplexen Ketten kann schwierig sein, insbesondere wenn viele Handler beteiligt sind. Die Logik zur Weiterleitung von Anfragen kann unübersichtlich werden.
}

\nachteil{Leistungseinbußen}
\note{
    Bei sehr langen Ketten oder vielen Handlern kann es zu Leistungseinbußen kommen, da jede Anfrage durch mehrere Handler verarbeitet werden muss. Dies kann die Reaktionszeiten erhöhen.
}

\nachteil{Unklare Verantwortlichkeiten}
\note{
    Wenn die Kette nicht gut dokumentiert oder organisiert ist, kann es schwierig sein zu verstehen, welcher Handler welche Aufgaben übernimmt und wie die Anfragen verarbeitet werden.
}

\nachteil{Gefahr unkontrollierter Weiterleitung}
\note{
    Wenn ein Handler die Anfrage nicht korrekt weiterleitet oder sich nicht an das erwartete Verhalten hält, kann dies zu unvorhersehbaren Ergebnissen oder Fehlern in der Kette führen.
}